\chapter{Introduction}

\section{Nightmares and Nightmare Disorder as a Clinical Problem}

\todo{Introduce nightmares as a clinically relevant disorder, including prevalence, diagnostic criteria, and psychological burden.}

\subsection{Prevalence and Clinical Burden}

\todo{Draft subsection.}

\subsection{Diagnostic Criteria of Nightmare Disorder}

\todo{Draft subsection.}

\subsection{Clinical Relevance for PTSD and Mental Health}

\todo{Draft subsection.}

\section{Imagery Rehearsal Therapy and the Treatment Gap}

\todo{Introduce IRT as an evidence-based treatment and motivate the gap between treatment recommendations and real-world access or awareness.}

\subsection{IRT as an Evidence-Based Treatment}

\todo{Draft subsection.}

\subsection{Limited Awareness and Access}

\todo{Draft subsection.}

\subsection{Motivation for Scalable Support}

\todo{Draft subsection.}

\section{AI-Assisted Mental Health Interventions}

\todo{Situate the thesis in the growing use of AI for mental-health support and motivate IRT as a structured but safety-sensitive use case.}

\subsection{Growing Use of AI for Mental Health Support}

\todo{Draft subsection.}

\subsection{Existing AI Mental Health Approaches}

\todo{Draft subsection.}

\subsection{Opportunity for AI-Assisted IRT}

\todo{Draft subsection.}

\section{The Safety Challenge in AI-Mediated IRT}

\todo{Develop the central problem: LLMs can produce fluent therapeutic-sounding responses that may still violate IRT principles.}

\subsection{General Limitations of LLMs}

\todo{Draft subsection.}

\subsection{Risks in Mental Health Contexts}

\todo{Draft subsection.}

\subsection{Risks in IRT Rescripting}

\todo{Draft subsection.}

\subsection{Need for an Explicit Clinical Safety Layer}

\todo{Draft subsection.}

\begin{figure}
  \centering
  \fbox{\parbox{0.85\linewidth}{\centering Figure placeholder: safety problem example in AI-assisted IRT.}}
  \caption{Example of the safety problem in AI-assisted IRT. A rescripting attempt that appears protective can reinforce avoidance when validated by an unconstrained LLM. The proposed safety layer identifies the maladaptive pattern and redirects the intervention toward adaptive mastery.}
  \label{fig:safety-problem-example}
\end{figure}

\section{Research Questions and Evaluation Strategy}

This thesis asks whether a neuro-symbolic, knowledge-graph-augmented safety layer can reduce clinically maladaptive content and improve IRT guideline adherence compared with an unconstrained baseline LLM.

\subsection{Main Research Question}

Can a neuro-symbolic, knowledge-graph-augmented safety layer reduce clinically maladaptive content and improve IRT guideline adherence compared with an unconstrained baseline LLM?

\subsection{RQ1: Detection and Classification Reliability}

Can the ontology-grounded classifier reliably identify adaptive and maladaptive rescripting patterns?

\subsection{RQ2: Classifier Design and Pipeline Components}

How do reasoning depth, preprocessing, and pipeline components affect safety classification performance?

\subsection{RQ3: Uncertainty Quantification}

Which uncertainty signal better predicts classifier errors: token log-probability or semantic consistency across stochastic runs?

\subsection{RQ4: System-Level Safety and Guideline Adherence}

Does the safety-augmented chatbot produce outputs rated as safer and more guideline-adherent than the unconstrained baseline?

\subsection{Success Criteria}

\todo{Define expected thresholds and interpretation criteria, including classifier accuracy, unsafe recall, uncertainty usefulness, ablation findings, and system-level expert ratings.}

\section{Contributions}

\todo{Summarize the thesis contributions: clinical safety ontology, benchmark dataset, neuro-symbolic safety pipeline, and methodological approach.}

\subsection{Clinical Safety Ontology for IRT Rescripting}

\todo{Draft subsection.}

\subsection{Benchmark Dataset for IRT Safety Classification}

\todo{Draft subsection.}

\subsection{Neuro-Symbolic Safety Critic Pipeline}

\todo{Draft subsection.}

\subsection{Methodological Contribution: Formalizing Protocol-Specific Clinical Safety}

\todo{Draft subsection.}

\section{Thesis Outline}

\todo{Briefly preview the remaining chapters.}
% !TEX root = ../main.tex
